\documentclass[11pt,a4paper]{article}
% Preambulo comun de los reportes de Tecnologias Celulares.
% Cada main.tex hace: \documentclass ... % Preambulo comun de los reportes de Tecnologias Celulares.
% Cada main.tex hace: \documentclass ... % Preambulo comun de los reportes de Tecnologias Celulares.
% Cada main.tex hace: \documentclass ... \input{../comun/preambulo.tex}
\usepackage[utf8]{inputenc}
\usepackage[T1]{fontenc}
\usepackage[spanish,mexico]{babel}
\usepackage{graphicx}
\usepackage[margin=2.5cm]{geometry}
\usepackage{amsmath}
\usepackage{amssymb}
\usepackage{booktabs}
\usepackage{tabularx}
\usepackage{array}
\usepackage{float}
\usepackage{setspace}
\usepackage{parskip}
\usepackage{fancyhdr}
\usepackage{textcomp}
\usepackage{xcolor}
\usepackage{caption}
\usepackage{fvextra}
\usepackage[hidelinks]{hyperref}

\onehalfspacing
\setlength{\parskip}{6pt}
\sloppy
\setlength{\emergencystretch}{3em}
\setlength{\headheight}{40pt}
\captionsetup{font=small,labelfont=bf}
\setlength{\tabcolsep}{6pt}

\pagestyle{fancy}
\fancyhf{}
\fancyhead[L]{\includegraphics[width=1cm]{/home/yovick/repos/UAZ-IEI/Assets/UIE.png}}
\fancyhead[R]{\small Tecnologías Celulares}
\fancyfoot[C]{\thepage}
\renewcommand{\headrulewidth}{0.4pt}

\DefineVerbatimEnvironment{codigo}{Verbatim}{%
  breaklines=true,breakanywhere=true,fontsize=\small,frame=single,
  framesep=3pt,rulecolor=\color{gray!50}}

% Figura centrada con pie: \figura[ancho]{archivo}{leyenda}
\newcommand{\figura}[3][0.85]{%
  \begin{center}
    \includegraphics[width=#1\linewidth]{#2}\\[2pt]
    {\small #3}
  \end{center}}

% Portada: \portada{numero}{titulo de la practica}
\newcommand{\portada}[2]{%
\begin{titlepage}
  \thispagestyle{empty}
  \begin{center}
    \vspace*{1cm}
    \includegraphics[width=4.5cm]{/home/yovick/repos/UAZ-IEI/Assets/uaz.png}
    \vspace{0.8cm}

    {\large Universidad Autónoma de Zacatecas}\\[4pt]
    {\large Unidad Académica de Ingeniería Eléctrica}\\[4pt]
    {\normalsize Ingeniería Electrónica Industrial con especialidad en Telecomunicaciones}

    \vspace{1.8cm}
    {\LARGE\bfseries Práctica #1}\\[10pt]
    {\Large\bfseries #2}
  \end{center}

  \vspace{2.2cm}
  \begin{center}
    \begin{tabular}{rl}
      \textbf{Materia:} & Tecnologías Celulares\\
      \textbf{Docente:} & Dr. Víktor Iván Rodríguez Abdalá\\
      \textbf{Alumno:}  & Arnold Yovick Ríos Zaldívar\\
      \textbf{Fecha:}   & Noviembre de 2026\\
    \end{tabular}
  \end{center}
  \vfill
\end{titlepage}}

\usepackage[utf8]{inputenc}
\usepackage[T1]{fontenc}
\usepackage[spanish,mexico]{babel}
\usepackage{graphicx}
\usepackage[margin=2.5cm]{geometry}
\usepackage{amsmath}
\usepackage{amssymb}
\usepackage{booktabs}
\usepackage{tabularx}
\usepackage{array}
\usepackage{float}
\usepackage{setspace}
\usepackage{parskip}
\usepackage{fancyhdr}
\usepackage{textcomp}
\usepackage{xcolor}
\usepackage{caption}
\usepackage{fvextra}
\usepackage[hidelinks]{hyperref}

\onehalfspacing
\setlength{\parskip}{6pt}
\sloppy
\setlength{\emergencystretch}{3em}
\setlength{\headheight}{40pt}
\captionsetup{font=small,labelfont=bf}
\setlength{\tabcolsep}{6pt}

\pagestyle{fancy}
\fancyhf{}
\fancyhead[L]{\includegraphics[width=1cm]{/home/yovick/repos/UAZ-IEI/Assets/UIE.png}}
\fancyhead[R]{\small Tecnologías Celulares}
\fancyfoot[C]{\thepage}
\renewcommand{\headrulewidth}{0.4pt}

\DefineVerbatimEnvironment{codigo}{Verbatim}{%
  breaklines=true,breakanywhere=true,fontsize=\small,frame=single,
  framesep=3pt,rulecolor=\color{gray!50}}

% Figura centrada con pie: \figura[ancho]{archivo}{leyenda}
\newcommand{\figura}[3][0.85]{%
  \begin{center}
    \includegraphics[width=#1\linewidth]{#2}\\[2pt]
    {\small #3}
  \end{center}}

% Portada: \portada{numero}{titulo de la practica}
\newcommand{\portada}[2]{%
\begin{titlepage}
  \thispagestyle{empty}
  \begin{center}
    \vspace*{1cm}
    \includegraphics[width=4.5cm]{/home/yovick/repos/UAZ-IEI/Assets/uaz.png}
    \vspace{0.8cm}

    {\large Universidad Autónoma de Zacatecas}\\[4pt]
    {\large Unidad Académica de Ingeniería Eléctrica}\\[4pt]
    {\normalsize Ingeniería Electrónica Industrial con especialidad en Telecomunicaciones}

    \vspace{1.8cm}
    {\LARGE\bfseries Práctica #1}\\[10pt]
    {\Large\bfseries #2}
  \end{center}

  \vspace{2.2cm}
  \begin{center}
    \begin{tabular}{rl}
      \textbf{Materia:} & Tecnologías Celulares\\
      \textbf{Docente:} & Dr. Víktor Iván Rodríguez Abdalá\\
      \textbf{Alumno:}  & Arnold Yovick Ríos Zaldívar\\
      \textbf{Fecha:}   & Noviembre de 2026\\
    \end{tabular}
  \end{center}
  \vfill
\end{titlepage}}

\usepackage[utf8]{inputenc}
\usepackage[T1]{fontenc}
\usepackage[spanish,mexico]{babel}
\usepackage{graphicx}
\usepackage[margin=2.5cm]{geometry}
\usepackage{amsmath}
\usepackage{amssymb}
\usepackage{booktabs}
\usepackage{tabularx}
\usepackage{array}
\usepackage{float}
\usepackage{setspace}
\usepackage{parskip}
\usepackage{fancyhdr}
\usepackage{textcomp}
\usepackage{xcolor}
\usepackage{caption}
\usepackage{fvextra}
\usepackage[hidelinks]{hyperref}

\onehalfspacing
\setlength{\parskip}{6pt}
\sloppy
\setlength{\emergencystretch}{3em}
\setlength{\headheight}{40pt}
\captionsetup{font=small,labelfont=bf}
\setlength{\tabcolsep}{6pt}

\pagestyle{fancy}
\fancyhf{}
\fancyhead[L]{\includegraphics[width=1cm]{/home/yovick/repos/UAZ-IEI/Assets/UIE.png}}
\fancyhead[R]{\small Tecnologías Celulares}
\fancyfoot[C]{\thepage}
\renewcommand{\headrulewidth}{0.4pt}

\DefineVerbatimEnvironment{codigo}{Verbatim}{%
  breaklines=true,breakanywhere=true,fontsize=\small,frame=single,
  framesep=3pt,rulecolor=\color{gray!50}}

% Figura centrada con pie: \figura[ancho]{archivo}{leyenda}
\newcommand{\figura}[3][0.85]{%
  \begin{center}
    \includegraphics[width=#1\linewidth]{#2}\\[2pt]
    {\small #3}
  \end{center}}

% Portada: \portada{numero}{titulo de la practica}
\newcommand{\portada}[2]{%
\begin{titlepage}
  \thispagestyle{empty}
  \begin{center}
    \vspace*{1cm}
    \includegraphics[width=4.5cm]{/home/yovick/repos/UAZ-IEI/Assets/uaz.png}
    \vspace{0.8cm}

    {\large Universidad Autónoma de Zacatecas}\\[4pt]
    {\large Unidad Académica de Ingeniería Eléctrica}\\[4pt]
    {\normalsize Ingeniería Electrónica Industrial con especialidad en Telecomunicaciones}

    \vspace{1.8cm}
    {\LARGE\bfseries Práctica #1}\\[10pt]
    {\Large\bfseries #2}
  \end{center}

  \vspace{2.2cm}
  \begin{center}
    \begin{tabular}{rl}
      \textbf{Materia:} & Tecnologías Celulares\\
      \textbf{Docente:} & Dr. Víktor Iván Rodríguez Abdalá\\
      \textbf{Alumno:}  & Arnold Yovick Ríos Zaldívar\\
      \textbf{Fecha:}   & Noviembre de 2026\\
    \end{tabular}
  \end{center}
  \vfill
\end{titlepage}}


\begin{document}

\portada{11}{Código espacio-tiempo de Alamouti para sistemas MIMO}

\textbf{Script:} \texttt{practica11.py} (Python con NumPy/SciPy)

\section{Objetivo}

Implementar un sistema de transmisión con el código espacio-tiempo de Alamouti (2 antenas
transmisoras y 1 receptora), evaluar su BER sobre un canal Rayleigh y comparar su desempeño
con un sistema SISO convencional para evidenciar la ganancia de diversidad.

\section{Fundamento teórico}

El código de Alamouti transmite dos símbolos $s_1$ y $s_2$ en dos intervalos consecutivos
desde dos antenas:
\begin{equation}
  \mathbf{S} =
  \begin{pmatrix}
    s_1 & -s_2^{*}\\
    s_2 & s_1^{*}
  \end{pmatrix},
\end{equation}
donde las filas son los intervalos de tiempo y las columnas las antenas. La matriz es
ortogonal:
\begin{equation}
  \mathbf{S}^{H}\mathbf{S} = \left(|s_1|^2+|s_2|^2\right)\mathbf{I}.
\end{equation}
Con coeficientes de canal $h_1$ y $h_2$ (constantes durante los dos intervalos) y ruido
$n_1,n_2$, las señales recibidas son
\begin{align}
  r_1 &= \frac{h_1 s_1 + h_2 s_2}{\sqrt{2}} + n_1,\\
  r_2 &= \frac{-h_1 s_2^{*} + h_2 s_1^{*}}{\sqrt{2}} + n_2.
\end{align}
El decodificador combina linealmente:
\begin{equation}
  \hat{s}_1 = \frac{h_1^{*}r_1 + h_2 r_2^{*}}{g},
  \qquad
  \hat{s}_2 = \frac{h_2^{*}r_1 - h_1 r_2^{*}}{g},
  \qquad
  g = \frac{|h_1|^2+|h_2|^2}{\sqrt{2}},
\end{equation}
y la interferencia entre $s_1$ y $s_2$ se cancela gracias a la ortogonalidad. El SNR
efectivo es $\gamma = (|h_1|^2+|h_2|^2)E_b/(2N_0)$, es decir, la suma de dos ramas
independientes: el sistema logra \textbf{diversidad de orden 2} manteniendo tasa unitaria
(dos símbolos en dos intervalos) y sin necesidad de conocer el canal en el transmisor.

Para comparar de forma justa, la potencia total radiada es la misma en ambos sistemas: el
SISO transmite $s=\pm1$ con energía unitaria y Alamouti transmite $s/\sqrt{2}$ por antena.
La BER teórica de Alamouti se obtiene promediando $Q(\sqrt{2\gamma})$ sobre la distribución
Gamma del SNR combinado.

\section{Desarrollo experimental}

\subsection{Parámetros}

\begin{tabularx}{\linewidth}{@{}*{2}{>{\raggedright\arraybackslash}X}@{}}
  \toprule
  Parámetro & Valor \\
  \midrule
  Antenas & 2 transmisoras, 1 receptora \\
  Modulación & BPSK \\
  Bits por punto & 200\,000 \\
  Canal & Rayleigh i.i.d., $\mathcal{CN}(0,1)$ \\
  Potencia total & unitaria (cada antena $1/\sqrt{2}$) \\
  $E_b/N_0$ & 0 a 20 dB en pasos de 2 dB \\
  \bottomrule
\end{tabularx}

\subsection{Procedimiento}

\begin{enumerate}
  \item Se generan bits y se agrupan en pares $(s_1,s_2)$.
  \item Se aplica la matriz de Alamouti y se transmiten los dos intervalos por las dos
        antenas.
  \item Se genera un canal Rayleigh independiente por par y se añade ruido complejo con
        $N_0=1/(E_b/N_0)$.
  \item Se decodifica con las combinaciones lineales, se normaliza por $g$ y se decide por
        el signo de la parte real.
  \item Se repite el experimento para un SISO con la misma potencia y se comparan las
        curvas.
\end{enumerate}

\section{Código}

\begin{codigo}
r1 = (h1 * s1 + h2 * s2) / np.sqrt(2) + ruido(rng, npares, n0)
r2 = (-h1 * np.conj(s2) + h2 * np.conj(s1)) / np.sqrt(2) + ruido(rng, npares, n0)

ganancia = (np.abs(h1) ** 2 + np.abs(h2) ** 2) / np.sqrt(2)
s1_hat = (np.conj(h1) * r1 + h2 * np.conj(r2)) / ganancia
s2_hat = (np.conj(h2) * r1 - h1 * np.conj(r2)) / ganancia
\end{codigo}

\section{Resultados}

\subsection{BER versus $E_b/N_0$}

\begin{tabular}{@{}ccc@{}}
  \toprule
  $E_b/N_0$ (dB) & BER SISO & BER Alamouti \\
  \midrule
  0  & 0.14705 & 0.11539 \\
  2  & 0.10820 & 0.07483 \\
  4  & 0.07564 & 0.04346 \\
  6  & 0.05228 & 0.02373 \\
  8  & 0.03553 & 0.01186 \\
  10 & 0.02387 & 0.00567 \\
  12 & 0.01518 & 0.00261 \\
  14 & 0.00970 & 0.00114 \\
  16 & 0.00646 & 0.00040 \\
  18 & 0.00388 & 0.00019 \\
  20 & 0.00260 & 0.00008 \\
  \bottomrule
\end{tabular}

\figura{figs/ber.png}{BER de SISO y Alamouti sobre canal Rayleigh, con sus curvas teóricas.}

La simulación coincide con la teoría en todo el rango. A 10 dB, Alamouti reduce la BER de
0.0239 a 0.0057 (factor 4.2). La curva de Alamouti cae con una pendiente mucho mayor: a
partir de 12 dB la diferencia se amplía hasta dos órdenes de magnitud a 20 dB.

\subsection{Ganancia de diversidad}

Para una BER de $10^{-2}$, el SISO requiere 13.86 dB y Alamouti 8.46 dB: una ganancia de
\textbf{5.4 dB}. Para $10^{-3}$, Alamouti alcanza el objetivo en 14.24 dB, mientras que el
SISO no llega a ese valor ni siquiera a 20 dB (su BER mínima simulada fue
$2.6\times10^{-3}$). La pendiente de la curva de Alamouti es de orden 2, frente al orden 1
del SISO, que es la firma de la diversidad espacial.

\subsection{Constelación detectada}

\figura[0.55]{figs/constelacion.png}{Símbolo $s_1$ detectado a $E_b/N_0=10$ dB.}

Los símbolos se concentran alrededor de $\pm1$ en el eje real. La dispersión es menor que
la de un SISO a la misma SNR porque el decodificador combina dos réplicas independientes
de la señal.

\section{Análisis de resultados}

\begin{itemize}
  \item \textbf{¿Qué ganancia proporciona Alamouti?} Una ganancia de diversidad de orden 2:
        la BER cae con la segunda potencia del SNR en lugar de la primera. En la práctica,
        son 5.4 dB a una BER de $10^{-2}$; a $10^{-3}$ el SISO ni siquiera alcanza el umbral en el rango simulado.
  \item \textbf{¿Por qué la BER disminuye?} Porque el receptor combina dos copias
        independientes de la señal (una por antena) y la probabilidad de que ambas sufran
        un desvanecimiento profundo al mismo tiempo es mucho menor que la de una sola.
  \item \textbf{¿Qué significa diversidad espacial?} Disponer de réplicas de la señal en
        antenas separadas, de modo que los desvanecimientos independientes de cada camino
        se promedian y reducen la probabilidad de error.
  \item \textbf{¿Qué ventajas ofrece frente a incrementar la potencia transmitida?} La
        diversidad mejora la BER de forma cuadrática con la SNR, mientras que aumentar la
        potencia solo da una mejora lineal (3 dB por cada duplicación). Además, la
        diversidad no consume más espectro ni más potencia total.
\end{itemize}

\section{Cuestionario}

\begin{enumerate}
  \item \textbf{¿Qué es un código espacio-tiempo?} Un esquema que introduce redundancia en
        el espacio (varias antenas) y en el tiempo (varios intervalos) para que el receptor
        obtenga diversidad y/o multiplexación sin sacrificar tasa.
  \item \textbf{¿Qué ventaja tiene el esquema de Alamouti?} Logra diversidad de orden 2 con
        tasa unitaria y un decodificador lineal muy simple, sin requerir conocimiento del
        canal en el transmisor.
  \item \textbf{¿Qué es la diversidad espacial?} La disponibilidad de varias rutas
        independientes entre transmisor y receptor gracias a antenas separadas; si una ruta
        se atenúa, las otras siguen llevando la señal.
  \item \textbf{¿Por qué la tasa del código Alamouti es unitaria?} Porque transmite dos
        símbolos en dos intervalos de tiempo; no repite información ni introduce redundancia
        temporal, solo la distribuye entre las antenas.
  \item \textbf{¿Qué aplicaciones celulares utilizan técnicas similares?} LTE usa
        transmisión con diversidad (SFBC, basada en Alamouti) para el canal de control y
        esquemas de diversidad de transmisión; 5G NR emplea diversidad espacial y
        precodificación con múltiples antenas.
\end{enumerate}

\section{Conclusiones}

Se implementó el código de Alamouti 2x1 y se comparó con un SISO sobre canal Rayleigh con
la misma potencia total. Alamouti presentó diversidad de orden 2: a 10 dB redujo la BER de
0.0239 a 0.0057, a $10^{-2}$ ganó 5.4 dB y a $10^{-3}$ alcanzó el objetivo en 14.24 dB,
valor que el SISO no logró en el rango simulado. El experimento confirma que la diversidad
espacial mejora la confiabilidad del enlace sin aumentar la potencia ni el ancho de banda.

\end{document}
